\documentclass{ximera}

% vier voorkeuren die je meteen zelf kan instellen:
\def\uitbr{0} % waarde 1 als je Uitbreiding in de marge wilt, 0 als je dat niet wilt 
\def\wsg{0} % waarde 1 als je verwijzingen naar Wiskunde Samen gevat² in de marge wilt, 0 als je dat niet wilt
\def\rectoverso{0} % waarde 1 als je het PDF-bestand recto-verso wilt laten afdrukken, 0 als je het recto wilt
\def\voetn{1} % waarde 1 als je de voetnoten wilt, 0 als je dat niet wilt 

% marges:
\usepackage[paper=a4paper,margin=3.4cm,marginparwidth=2cm]{geometry}

% packages algemeen:
\pdfOnly{
\usepackage[english,dutch]{babel}
}
\usepackage{amsfonts}
\usepackage{amsmath} %  o.a. \eqref en \text, omgeving equation*
\usepackage{amssymb} % o.a. \nmid
\usepackage{amsthm} % o.a. omgeving proof
\usepackage{graphicx} % o.a. figuren
\usepackage{xcolor} % \color
\usepackage[disable]{todonotes} % \todo 
\usepackage{comment} % omgeving comment
\usepackage{xcolor} % kleuren 
\usepackage[skip=7pt plus1pt, indent=0pt]{parskip} % skip = verticale ruimte tussen twee alinea's, indent = insprong bij nieuwe alinea
\usepackage{multicol} % kolommen
\usepackage{xfrac} % breuken

\usepackage{siunitx} % SI eenheden
\usepackage{eso-pic} % achtergrond voorpagina
\usepackage{tabularx} % kolommen voor rekenmachine
\usepackage{rotating} % voor omgeving turn

% voor figuren met PSTricks:
\usepackage{pstricks} 
\usepackage{pstricks-add}
\usepackage{pst-plot}
\usepackage{pst-node}
\usepackage{pst-coil}
\usepackage{auto-pst-pdf}

% voor figuren met TikZ:
\usepackage{tikz} 
\usepackage{tkz-euclide} 
\usepackage{pgfplots} 
\usetikzlibrary{calc,intersections,through,backgrounds,patterns} 
\pgfplotsset{compat=newest}
\usepgfplotslibrary{fillbetween,colormaps}
\usepackage{import}
\usepackage{tikz-3dplot}

% zelfde parskip in minipage:
\makeatletter
\setlength{\parskip}{\medskipamount}
\newcommand{\@minipagerestore}{\setlength{\parskip}{\medskipamount}}
\makeatother

% voor veranderen van de naam Bibliografie naar Referentielijst:
\pdfOnly{
\addto\captionsdutch{\renewcommand{\bibname}{Referentielijst}}

% voor trefwoordenlijst:
\usepackage{imakeidx}
\makeindex[title=Trefwoordenlijst,program=makeindex,options=-s index.tex,columns=2,intoc=true]
}
% voor nummeren van vergelijkingen:
%WIM%\numberwithin{equation}{chapter} % als je dit desactiveert dan worden vergelijkingen in bijvoorbeeld hoofdstuk 3 genummerd als (1), (2) etc. in plaats van (3.1), (3.2) etc.

% voor hyperlinks en bladwijzers in PDF-bestand:
%WIM%\usepackage[bookmarksopen,bookmarksopenlevel=0,hypertexnames=false,pdfa,bookmarksnumbered]{hyperref} 
%WIM%\usepackage{bookmark} 

% voor hyperlinks: 
\hypersetup{pdfborder={0 0 0}, pdfstartpage=1,linkbordercolor={1 0 0}}%pdfborder={0 0 0} is geen rand rond hyperlinks, pdfborder={1 1 1} wel

% korter commando voor displaystyle, om bijvoorbeeld breuken groter te drukken in doorlopende tekst: 
\newcommand{\D}{\displaystyle}

% (veel)gebruikte verzamelingen:
\newcommand\NN{\mathbb{N}} % verzameling van de natuurlijke getallen
\newcommand\QQ{\mathbb{Q}} % verzameling van de rationale getallen
\newcommand\RR{\mathbb{R}} % verzameling van de re\"ele getallen
\newcommand\ZZ{\mathbb{Z}} % verzameling van de gehele getallen

% (veel)gebruikte operatoren:
\def\co{\operatorname{co}} % coördinaat van een punt
\def\ggd{\operatorname{ggd}} % positieve grootste gemene deler van twee gehele getallen niet beide nul
\def\gr{\operatorname{gr}} % graad van een veelterm

% voor kleur grafieken (donkergroen):
\definecolor{graf}{RGB}{0,100,0} 

% voor lijsten: 
% \usepackage{enumerate}
%WIM%\usepackage[shortlabels]{enumitem} 
\setlist{topsep=0em, itemsep=-0.15em}

% voor small bullet:
\newcommand\sbullet[1][.5]{\mathbin{\vcenter{\hbox{\scalebox{#1}{$\bullet$}}}}}

% voor meer verticale ruimte tussen vergelijkingen in align
\addtolength{\jot}{0.1cm}

% voor meervoudige voetnoten:
\usepackage{fnpct}

% voor het onderdrukken van voetnoten als optie:
\usepackage{letltxmacro}
\LetLtxMacro\Oldfootnote\footnote
\newcommand{\EnableFootNotes}{%
  \LetLtxMacro\footnote\Oldfootnote%
}
\newcommand{\DisableFootNotes}{%
  \renewcommand{\footnote}[2][]{\relax}
}

% voor verticale lijn in de marge (uitbreiding):
\usepackage[framemethod=default]{mdframed}
\usepackage{marginnote}
\reversemarginpar
\ifthenelse{\uitbr < 1}{\definecolor{rood}{RGB}{255,255,255}}{\definecolor{rood}{RGB}{254,64,64}} %HEX: #fe4040
\mdfdefinestyle{uitbreiding}{%
    topline=false,
    rightline=false,
    bottomline=false,
    leftline=true,
    linecolor=rood,
    linewidth=5pt,
    rightmargin=0pt,
    skipabove=10pt,% ipv 3
    skipbelow=0pt,
    leftmargin=-25pt,
    innerleftmargin=20pt,
    innerrightmargin=0pt,
    innertopmargin=0pt,
    innerbottommargin=0pt%,
%	needspace=30pt %minimumhoogte vooraleer lijn wordt gesplitst
    }
\newenvironment{Uitbreiding}{
    \marginpar{
        \center
		\vspace{0.1cm}
		\vspace{7pt}
		\rotatebox{90}{\color{rood}\Large \bf Uitbreiding}
	}
    \begin{mdframed}[style=uitbreiding]
    }{\vspace{-0.05cm}
    \end{mdframed}
}


% omgevingen voor lemma, definitie, voorbeeld etc.:
\newtheoremstyle{mystyle}
    {0em} % Space above
    {0em} % Space below
    {\itshape} % Body font
    {} % Indent amount
    {\bfseries} % Theorem head font
    {.} % Punctuation after theorem head
    {.5em} % Space after theorem head
    {} % Theorem head spec (can be left empty, meaning `normal')
\theoremstyle{mystyle}
%WIM%\newtheorem{lemma}{Lemma}[chapter] % als je [chapter] desactiveert dan Voorbeeld 3 in plaats van Voorbeeld 5.3, als je [chapter] vervangt door [section] dan Voorbeeld 5.2.3 in plaats van Voorbeeld 5.3 
%\theoremstyle{definition} % als je dit activeert dan is wat in de omgeving staat niet cursief gedrukt
\newtheorem{oefening}[lemma]{Oefening}
\newtheorem{definitie}[lemma]{Definitie} 
\newtheorem{voorbeeld}[lemma]{Voorbeeld}
\newtheorem{eigenschap}[lemma]{Eigenschap} 
\newtheorem{stelling}[lemma]{Stelling} 
\newtheorem{gevolg}[lemma]{Gevolg} 
\newtheorem{afspraak}[lemma]{Afspraak} 
\newtheorem{werkwijze}[lemma]{Werkwijze}

% omgeving proof, aangepaste ruimte: 
\makeatletter
\renewenvironment{proof}[1][\proofname]{\par
  \vspace{-\topsep}% remove the space after the theorem
  \pushQED{\qed}%
  \normalfont
  \topsep0pt \partopsep0pt % no space before
  \trivlist
  \item[\hskip\labelsep
        \itshape
    #1\@addpunct{.}]\ignorespaces
}{%
  \popQED\endtrivlist\@endpefalse
  \addvspace{0pt plus 0pt} % no space after
}
\makeatother

% kaderstijlen uit SOHO Wiskunde Plantyn:
\colorlet{steunkleur}{black}
\colorlet{steunkleurlicht}{steunkleur!30!white}
\colorlet{steunkleurkader}{steunkleur!7!white}
\colorlet{steunkleurkaderlicht}{steunkleur!2!white}

\mdfdefinestyle{kaderstijl_vol_licht}
{skipabove=6pt, 
skipbelow=6pt, 
backgroundcolor=steunkleurkader,
linecolor=steunkleurlicht,
linewidth = 0.4pt, 
topline=true,
bottomline=true, 
rightline=true,
innerleftmargin=5pt,
innerrightmargin=5pt,
innertopmargin=5pt,
leftmargin=0cm,
rightmargin=0cm,
innerbottommargin=5pt,
needspace=30pt % minimumhoogte voor splitsen kader
}
\surroundwithmdframed[style=kaderstijl_vol_licht]{definitie}
\surroundwithmdframed[style=kaderstijl_vol_licht]{stelling}
\surroundwithmdframed[style=kaderstijl_vol_licht]{lemma}
\surroundwithmdframed[style=kaderstijl_vol_licht]{eigenschap}
\surroundwithmdframed[style=kaderstijl_vol_licht]{gevolg}

% voor omgevingen voor oefening en antwoord:
\newenvironment{Oefening}{%
    \begin{enumerate}[ 
    series=Oef,
    resume=Oef,
    leftmargin=1.78em,
    label={\bfseries\arabic*.},
    ref=\arabic*
    ]
    \item %
    }{%
    \end{enumerate}
}
\newenvironment{Antwoord}{%
    \begin{enumerate}[%
    series=Antw,
    resume=Antw,
    leftmargin=1.78em,
    label={\bfseries\arabic*.},
    ref=\arabic*
    ]
    \item %
    }{%
    \end{enumerate}
}

% als de omgeving Uitbreiding start met een kaderomgeving (definitie, stelling, eigenschap, lemma of gevolg) dan moet wat extra verticale ruimte voorzien worden, met het commando \uitbreidingstartmetkader:
\def\uitbreidingstartmetkader{\mbox{}\vspace{-0.205cm}}

% % voor schema van de staartdeling:
% \usepackage{stackengine}
% \setstackgap{S}{5pt}
% \stackMath\def\stackalignment{r}
% \newcommand{\myRule}[3][white]{\textcolor{#1}{\rule{#2}{#3}}}
% \let\ph\phantom % enkel voor tekst
% \newcommand{\mph}[1]{% enkel voor math mode
%     \mathcolor{white}{#1}%
% }
% \def\staartmin{\rule{0.25cm}{0.1mm}\myRule{0.3cm}{0.1mm}}
% \def\staartphmin{\myRule{0.25cm}{0.1mm}\myRule{0.3cm}{0.1mm}}
% \newcommand{\staartstreep}[1]{\rule{\widthof{$#1$}}{0.1mm}}
% \newcommand{\staartphstreep}[1]{\myRule{\widthof{$#1$}}{0.1mm}}

% voor kolommen met schema van Horner:
\newcommand{\kolbreed}{1.0cm}
\newcolumntype{H}{>{\centering\arraybackslash$} p{\kolbreed} <{$}}

% voor nieuw commando utikzdashed: onderlijn (zoals underline) maar dan met stippellijn:
\tikzset{
    cheating dash/.code args={on #1 off #2 ends #3}{%
        \csname tikz@addoption\endcsname{%
            \pgfgetpath\currentpath%
            \pgfprocessround{\currentpath}{\currentpath}%
            \csname pgf@decorate@parsesoftpath\endcsname{\currentpath}{\currentpath}%
            \pgfmathparse{max(#1-#3,0)}\let\dashphase=\pgfmathresult%
            \pgfmathparse{\csname pgf@decorate@totalpathlength\endcsname-#1+2*\dashphase}\let\rest=\pgfmathresult%
            \pgfmathparse{#1+#2}\let\onoff=\pgfmathresult%
            \pgfmathparse{max(floor(\rest/\onoff), 1)}\let\nfullonoff=\pgfmathresult%
            \pgfmathparse{max((\rest-\onoff*\nfullonoff)/\nfullonoff+#2, #2)}\let\offexpand=\pgfmathresult%
            \pgfsetdash{{#1}{\offexpand}}{\dashphase pt}}%
    },
    cheating dash per segment/.style args={on #1 off #2 ends #3}{
        /utils/exec=\csname tikz@options\endcsname,%
        decoration={show path construction,
            lineto code={\draw [cheating dash=on #1 off #2 ends #3] (\tikzinputsegmentfirst) -- (\tikzinputsegmentlast);},
            curveto code={\draw [cheating dash=on #1 off #2 ends #3] (\tikzinputsegmentfirst) .. controls (\tikzinputsegmentsupporta) and (\tikzinputsegmentsupportb) .. (\tikzinputsegmentlast);},
            closepath code={\draw [cheating dash=on #1 off #2 ends #3] (\tikzinputsegmentfirst) -- (\tikzinputsegmentlast);}
        },
        decorate,
    },
}
\newcommand{\utikzdash}[1]{%
    \tikz[baseline=(todotted.base)]{
        \node[inner sep=0pt,outer sep=1.5pt] (todotted) {#1};
        \draw[cheating dash per segment=on 2pt off 2pt ends 2pt, line width=0.4pt] (todotted.south west) -- (todotted.south east);
    }%
}%

\providecommand{\xmemph}[1]{\textit{#1}}

% voor nieuw commando underdashed: onderlijn met stippellijnen en ook het woord cursief zetten:
\newcommand{\underdashed}[1]{%
    {\em\utikzdash{\!#1\!}}%
}

% voor icoon TI-84 Plus met verwijzing naar filmpje:
\newcommand{\grmlink}{\raisebox{0cm}{\includegraphics[width=1cm]{TI84Plus-icoon}}}
\newcommand{\grmref}[1]{%
	\reversemarginpar%
	\marginpar{
		\vspace{-0.4cm}%
		\htmladdnormallink{\grmlink}{#1}
	}	
}

% GRM knoppen:
\newcommand{\GRM}[1]{\fbox{\rule[0mm]{0cm}{0.215cm}\textup{\texttt{#1}}}}
\newcommand{\wedgetext}{{\raisebox{0.02cm}{\begin{turn}{90}>\end{turn}}}}
\newcommand{\veetext}{\raisebox{0.2cm}{\begin{turn}{-90}>\end{turn}}}

% voor kolommen met GRM screens:
\newlength{\widthallscreens}
\newlength{\widthscreens}
\newlength{\spaceleftscreen}
\newlength{\spacerightscreen}
\newlength{\spacebetweenscreens}

\newcommand{\setscreens}{
	\setlength{\spaceleftscreen}{2pt}
	\setlength{\spacerightscreen}{2pt}
	\setlength{\spacebetweenscreens}{8pt}
	\addtolength{\linewidth}{-28pt} % 3*8pt tussen vier screens en 2 pt links en 2 pt rechts
	\setlength{\widthallscreens}{\linewidth}
	\setlength{\widthscreens}{0.25\widthallscreens}
	\addtolength{\linewidth}{28pt}
	\newcolumntype{G}{p{\widthscreens}}
	\newcolumntype{s}{p{\spacebetweenscreens}}
	\newcolumntype{L}{p{\spaceleftscreen}}
	\newcolumntype{R}{p{\spacerightscreen}}
	\setlength{\tabcolsep}{0pt}
}

% voor icoon met verwijzing naar Wiskunde Samen gevat² in de marge:
% \newcommand{\wsglink}{\raisebox{0cm}{\includegraphics[width=1cm]{wsglogo}}}
\newcommand{\wsglink}{\raisebox{0cm}{LOGO}}
\newcommand{\htmladdnormallink}[2]{\href{#2}{#1}}
\newcounter{pagnrwsg}
\newcommand{\wsgref}[3]{% 
    % #1 woord dat onderstippeld wordt
    % #2 pagina van Wiskunde Samen gevat² waar je op terecht komt als je op de link klikt  
    % #3 afgedrukt op het logo van Wiskunde Samen gevat² in de marge: één of meerdere paginanummers
    \ifthenelse{\wsg < 1}{#1}{%
        \underdashed{#1}%
        {\setcounter{pagnrwsg}{#2}}%
        {\addtocounter{pagnrwsg}{14}}%
        \reversemarginpar%
        \marginpar{\vspace{-0.6cm}%
           \htmladdnormallink{\wsglink}{https://online.fliphtml5.com/sanky/laea/\#p=\arabic{pagnrwsg}}%
            \raisebox{0.41cm}[0cm][0cm]{%
                \hspace{-1.5cm}\makebox[2cm][c]{\colorbox{white}{\texttt{\footnotesize{#3}}}}%
            }%
        }%
    }%
}

% voor invoegen blanco pagina bij de optie recto-verso:
\def\blancobijrectoverso{
	\ifthenelse{\rectoverso < 1}{\clearpage}{
	\clearpage
	\thispagestyle{empty}
	\mbox{}
	\clearpage
	}
}

% aantal pagina's van het bestand:
\ifthenelse{\rectoverso < 1}{\def\totpag{57}}{\def\totpag{66}}

% documenteigenschappen:
% \usepackage{hyperxmp}
% \hypersetup{
% pdftitle={Open Source Wiskunde Aan zet: Veeltermen}, 
% pdfnumpages={\totpag},
% pdfauthor={Koen De Naeghel},
% pdflang={nl},
% pdfkeywords={wiskunde, open source, wiskunde aan zet, veeltermen, secundair onderwijs, tweede graad, doorstroomfinaliteit},
% pdfsubject={Veeltermen},
% pdfcopyright={\unichar{"24B8} 2024 Koen De Naeghel},
% pdfdate={17 december 2024},
% pdfapart=1,
% pdfstartview=}

% voor het benoemen van titel, auteur en datum:
% \title{Veeltermen}
% \author{Auteur: Koen De Naeghel}
% \date{\today}

\newcommand\BackgroundPic{
    \put(-260,-125){
    \parbox[b][\paperheight]{\paperwidth}{%
    \vfill
    \centering
    \includegraphics[height=\paperwidth, keepaspectratio]{WaZlogo}%
    \vfill
}}}

\addPrintStyle{..}
\begin{document}
	\author{Wiskunde Op Maat}
	\xmtitle{Kwadratische vergelijkingen oplossen}{}

Met deze pagina kan je het oplossen van tweedegraadsvergelijkingen in één onbekende met voorschrift \(ax2+bx+c=0\) inoefenen. 
Onder 'download' kan je de pdf met of zonder antwoorden downloaden.

\newcommand{\choicetwee}{{\wordChoice{\choice[correct]{twee oplossingen}\choice{één oplossing}\choice{geen oplossingen}}}}
\newcommand{\choiceeen}{{\wordChoice{\choice{twee oplossingen}\choice[correct]{één oplossing}\choice{geen oplossingen}}}}
\newcommand{\choicegeen}{{\wordChoice{\choice{twee oplossingen}\choice{één oplossing}\choice[correct]{geen oplossingen}}}}

\newcommand{\choicepositief}{{\wordChoice{\choice[correct]{positief}\choice{nul}\choice{negatief}}}}
\newcommand{\choicenul}{{\wordChoice{\choice{positief}\choice[correct]{nul}\choice{negatief}}}}
\newcommand{\choicenegatief}{{\wordChoice{\choice{positief}\choice{nul}\choice[correct]{negatief}}}}
 
% positieve discriminant 
\begin{exercise}

Bepaal de oplossingen van de vierkantsvergelijking \(x^2 - 4x + 3 = 0\).   

\begin{question}
De discrimant is \choicepositief. 

\begin{feedback}
    Het algemene voorschrift van een tweede graadsvergelijking wordt gegeven door  \(ax^2 + bx + c = 0\). 
    Om een tweedegraadsvergelijking op te lossen bereken je eerst de discriminant \(\Delta = b^2 - 4ac\). 
    Het aantal oplossingen wordt bepaald door het teken van deze discriminant. 
    
\end{feedback}
\end{question}

\begin{question}
    Deze vergelijking heeft \choicetwee in de reële getallen. 
    \begin{feedback}
        Indien de discriminant \(\Delta = b^2 - 4ac\) positief is, zijn er twee oplossingen in de reële getallen. 
    \end{feedback}
\end{question}

\begin{question}
    Bepaal de wortels. 
    
    \begin{hint}
        
        Voor een tweedegraadsvergelijking  \(ax^2 - bx + c = 0\) waarbij de discriminant \( \Delta = b^2 - 4ac \) positief is hebben we altijd twee oplossing gegeven door: 
        
        \[
            x_{1} = \frac{-b + \sqrt{\Delta}}{2a}  \text{ en }  x_{2} = \frac{-b - \sqrt{\Delta}}{2a}
            \]
            
            Waarbij:
            \begin{itemize}
                \item \( a \) de coëfficiënt van \( x^2 \) is,
                \item \( b \) de coëfficiënt van \( x \) is,
                \item \( c \) de constante term is,
            \end{itemize}
            
    \end{hint}
    

\end{question}

\begin{oplossing}
            
In het algemeen wordt een tweedegraadsvergelijking in één onbekende gegeven door het voorschrift \(ax^2 - bx + c = 0\). 
De tweedegraadsvergelijking \(x^2 - 4x + 3 = 0\) heeft als coëfficiënten gelijk aan \(a = 1, b = -4 \text{ en } c = 3\). 
We zoeken de getalwaarden voor \(x\) waarvoor deze vergelijking gelijk is aan nul. 
Hiervoor kan de formule met discrimant gebruikt worden. 

\vspace{5mm}

De discrimant \( \Delta = b^2 - 4ac \) is gelijk aan \(\Delta = (-4)^2 -4\cdot 1 \cdot 3\). De vierkanswortel wordt dan gegeven door \(\sqrt{\Delta} = \sqrt{4} = 2\). Er zijn dus twee oplossingen omdat de discriminant positief is. Deze oplossingen worden gegeven door: 

\[
x_{1} = \frac{-b + \sqrt{\Delta}}{2a}  \text{ en }  x_{2} = \frac{-b - \sqrt{\Delta}}{2a}
\]

Invullen levert: 

\[
x_{1} = \frac{-(-4) + \sqrt{4}}{2\cdot 1}  = 3 
\]


\[
x_{2} = \frac{-(-4) - \sqrt{4}}{2\cdot 1}  = 1
\]
 

Als je de vergelijking beschouwt als een functievoorschrift, zijn deze oplossingen de snijpunten van de parabool \( f(x) = x^2 - 4x + 3 \) met de \(x\)-as. Aangezien de discriminant positief is, zal de parabool de \(x\)-as 2 keer snijden. 

\begin{image}
\begin{tikzpicture}[scale=1]

    % Draw axes
    \draw[->] (-1,0) -- (5,0) node[right] {\(x\)};
    \draw[->] (0,-1) -- (0,4) node[above] {\(y\)};
    
    % Draw the parabola y = x^2 - 4x + 3
    \draw[domain=-0.2:4.2, samples=100, smooth, very thick, blue] 
        plot (\x, {(\x)^2 - 4*(\x) + 3}); 
        
    \draw (2, 3) node[blue] {\( f(x) = x^2 - 4x + 3 \)};

    % Roots (zeros) of the parabola
    \filldraw[red] (1,0) circle (2pt) node[below] {\(x=1\)};
    \filldraw[red] (3,0) circle (2pt) node[below] {\(x=3\)};

\end{tikzpicture}
\end{image}

Met de discrimantformule bereken je dus deze snijpunten met de \(x\)-as zonder grafiek van de functie . 
\end{oplossing}
  
    
\end{exercise}


% nul discriminant 
\begin{exercise}

    Bepaal de oplossingen van de vierkantsvergelijking \( x^2 - 4x + 4 = 0 \) .

    
    \begin{question}
    De discriminant is \choicenul. 
    \begin{feedback}
        Het algemene voorschrift van een tweede graadsvergelijking wordt gegeven door  \(ax^2 + bx + c = 0\). 
        Om een tweedegraadsvergelijking op te lossen bereken je eerst de discriminant \(\Delta = b^2 - 4ac\). 
        Het aantal oplossingen wordt bepaald door het teken van deze discriminant. 
    \end{feedback}
    \end{question}
    
    \begin{question}
    Deze vergelijking heeft \choiceeen in de reële getallen. 
    \begin{feedback}
        Omdat de discriminant nul is hebben we een dubbel nulpunt. 
    \end{feedback}
    \end{question}
    
    \begin{question}
    Bepaal de wortels.
    \begin{hint}
        Voor een tweedegraadsvergelijking \(ax^2 - bx + c = 0\) met \(\Delta = 0\) geldt:
        \[
        x = \frac{-b}{2a}.
        \]
    \end{hint}
    \end{question}
    
    \begin{oplossing}
    In het algemeen wordt een tweedegraadsvergelijking gegeven door het voorschrift \( ax^2 - bx + c = 0 \) 
    Voor de vergelijking 
    \[
    x^2 - 4x + 4 = 0
    \]
    hebben we:
    \[
    a = 1,\quad b = -4,\quad c = 4.
    \]
    
    \vspace{3mm}
    
    Bereken de discriminant:
    \[
    \Delta = b^2 - 4ac = (-4)^2 - 4\cdot1\cdot4 = 16 - 16 = 0.
    \]
    
    Aangezien \(\Delta = 0\) is er één oplossing (dubbele wortel), gegeven door:
    \[
    x = \frac{-b}{2a} = \frac{-(-4)}{2\cdot1} = \frac{4}{2} = 2.
    \]
    
    \vspace{3mm}
    
    Als functievoorschrift is 
    \[
    f(x)=x^2 - 4x + 4
    \]
    een parabool die de \(x\)-as raakt in het punt \( (2,0) \).
    
    \begin{image}
    \begin{tikzpicture}[scale=1]
        % Teken de assen
        \draw[->] (-1,0) -- (5,0) node[right] {\(x\)};
        \draw[->] (0,-1) -- (0,4) node[above] {\(y\)};
        
        % Teken de parabool f(x) = x^2 - 4x + 4
        \draw[domain=-0.2:4.2, samples=100, smooth, very thick, blue] 
            plot (\x, {(\x)^2 - 4*(\x) + 4});
            
        \draw (2, 3) node[blue] {\( f(x) = x^2 - 4x + 4 \)};
        
        % Markeer de dubbele wortel
        \filldraw[red] (2,0) circle (2pt) node[below] {\(x=2\)};
        
    \end{tikzpicture}
    \end{image}
    
    \end{oplossing}
      
    
    % \begin{expandable}
    %     \begin{remark}
    %         Aan het functievoorschrift was de dubbele wortel af te lezen als merkwaardig prodcuct aangezien \( f(x) = x^2 - 4x + 4  = (x -2)^2\)
    %     \end{remark}
        
    % \end{expandable}    
    
\end{exercise}


% negatieve disciminant 
\begin{exercise}

    Bepaal de oplossingen van de vierkantsvergelijking \( x^2 - 4x + 5 = 0 \). 
    
    
    \begin{question}
    De discriminant is \choicenegatief.
    \begin{feedback}
        Voor een algemene tweedegraadsvergelijking \(ax^2 - bx + c = 0\) bereken je de discriminant als \(\Delta = b^2 - 4ac\). 
        Hier geldt: 
        \[
        \Delta = (-4)^2 - 4\cdot1\cdot5 = 16 - 20 = -4.
        \]
        Omdat \(\Delta < 0\), is de discriminant negatief.
    \end{feedback}
    \end{question}
    
    \begin{question}
    Deze vergelijking heeft \choicegeen  in de reële getallen.
    \begin{feedback}
        Wanneer de discriminant negatief is, zijn er oplossingen in de reële getallen. 
        In de algemene formule staat de discriminant onder een wortel. 
        Er is geen reëel getal waarvan het kwadraat negatief is en je kan dus geen oplossingen bepalen: 
        \[
        x_{1} = \frac{-b + \sqrt{\Delta}}{2a}  \text{ en }  x_{2} = \frac{-b - \sqrt{\Delta}}{2a}
        \]

    \end{feedback}
    \end{question}
    
    \begin{oplossing}
    Voor de vergelijking \( x^2 - 4x + 5 = 0 \)
    
    hebben we:
    \[
    a = 1,\quad b = -4,\quad c = 5.
    \]
    
    De discriminant wordt berekend als:
    \[
    \Delta = b^2 - 4ac = (-4)^2 - 4\cdot1\cdot5 = 16 - 20 = -4.
    \]
    
    Aangezien \(\Delta < 0\) zijn er geen reële oplossingen. 

    De grafiek van de parabool ligt volledig boven de \(x\)-as. Er zijn geen snijpunten. 
    \begin{image}
        \begin{tikzpicture}[scale=1]
            % Teken de assen
            \draw[->] (-1,0) -- (5,0) node[right] {\(x\)};
            \draw[->] (0,-1) -- (0,6) node[above] {\(y\)};
            
            % Teken de parabool f(x) = x^2 - 4x + 5
            \draw[domain=-0.2:4.2, samples=100, smooth, very thick, blue] 
                plot (\x, {(\x)^2 - 4*(\x) + 5});
                
            \draw (2, 5) node[blue] {\( f(x) = x^2 - 4x + 5 \)};
            
        \end{tikzpicture}
    \end{image}

    
    \end{oplossing}
    
\end{exercise}
    

% \begin{expandable}{remark}{De constante factor \(c\)}
% \begin{remark}
    
    
%     In het algemeen voorschrift \( ax^2+bx+c=0 \) geeft de constante term \(c\) een verticale verschuiving weer. 
%     De waarde van \(c\) heeft dus invloed op het aantal snijpunten met de \(x\)-as en hiermee kan je inzien waar die \(c\) ook invloed heeft op het tegen van de discriminant \(\Delta = b^2 - 4ac\). 

% \usepackage{tikz}

%     \begin{image}

%         \begin{scope}[xshift=0cm]
%         \begin{tikzpicture}[scale=1]

%             % Draw axes
%             \draw[->] (-1,0) -- (5,0) node[right] {\(x\)};
%             \draw[->] (0,-1) -- (0,4) node[above] {\(y\)};
            
%             % Draw the parabola y = x^2 - 4x + 3
%             \draw[domain=-0.2:4.2, samples=100, smooth, very thick, blue] 
%                 plot (\x, {(\x)^2 - 4*(\x) + 3}); 
                
%             \draw (2, 3) node[blue] {\( f(x) = x^2 - 4x + 3 \)};
        
%             % Roots (zeros) of the parabola
%             \filldraw[red] (1,0) circle (2pt) node[below] {\(x=1\)};
%             \filldraw[red] (3,0) circle (2pt) node[below] {\(x=3\)};
        
%         \end{tikzpicture}
%         \end{scope}

%         \begin{scope}[xshift=3cm]
%         \begin{tikzpicture}[scale=1]
%             % Teken de assen
%             \draw[->] (-1,0) -- (5,0) node[right] {\(x\)};
%             \draw[->] (0,-1) -- (0,4) node[above] {\(y\)};
            
%             % Teken de parabool f(x) = x^2 - 4x + 4
%             \draw[domain=-0.2:4.2, samples=100, smooth, very thick, blue] 
%                 plot (\x, {(\x)^2 - 4*(\x) + 4});
                
%             \draw (2, 3) node[blue] {\( f(x) = x^2 - 4x + 4 \)};
            
%             % Markeer de dubbele wortel
%             \filldraw[red] (2,0) circle (2pt) node[below] {\(x=2\)};
            
%         \end{tikzpicture}
%         \end{scope}


%         \begin{scope}[xshift=6cm]    
%         \begin{tikzpicture}[scale=1]
%             % Teken de assen
%             \draw[->] (-1,0) -- (5,0) node[right] {\(x\)};
%             \draw[->] (0,-1) -- (0,6) node[above] {\(y\)};
            
%             % Teken de parabool f(x) = x^2 - 4x + 5
%             \draw[domain=-0.2:4.2, samples=100, smooth, very thick, blue] 
%                 plot (\x, {(\x)^2 - 4*(\x) + 5});
                
%             \draw (2, 5) node[blue] {\( f(x) = x^2 - 4x + 5 \)};
            
%         \end{tikzpicture}
%         \end{scope}
        
%     \end{image}
    
    
%     Voor de geïntresseerd leerlingen: de grafiek boven de \(x\)-as met een negatieve discrimant in de vierkantswortel heeft geen oplossingen in de reële getallen. 
%     In de complexe getallen hebben wiskundigen echter wel oplossingen gevonden...

%     \begin{expandable}{youtube}{A very gentle introduction to complex numbers.}
%         \youtube{https://youtu.be/f079K1f2WQk}
%     \end{expandable}   

    
% \end{remark}
% \end{expandable}


\begin{exercise} Bepaal de oplossingen van volgende tweedegraadsvergelijkingen. 
  
    \begin{question} De wortels van \( x^2 - 4x + 3    = 0 \) \begin{uitkomst} Er zijn twee reële oplossingen \( 1 \) en  \( 3 \)          \end{uitkomst}\end{question}
    \begin{question} De wortels van \( -2x^2 + 4x + 8  = 0 \) \begin{uitkomst} Er zijn twee reële oplossingen \( -2 \) en  \(4 \)          \end{uitkomst}\end{question}
    \begin{question} De wortels van \( 2x^2 - 8x + 16  = 0 \) \begin{uitkomst} De dubbele wortel is gelijk aan \( 2 \)                     \end{uitkomst}\end{question}
    \begin{question} De wortels van \( 2x^2 + 5x + 2   = 0 \) \begin{uitkomst} Er zijn twee reële oplossingen \( -\frac{1}{2}\) en \(-2 \) \end{uitkomst}\end{question}
    \begin{question} De wortels van \( 3x^2 - 6x + 3   = 0 \) \begin{uitkomst} De dubbele wortel is gelijk aan \( 1   \)                   \end{uitkomst}\end{question}

\end{exercise}


\begin{exercise} Bepaal de oplossingen van volgende tweedegraadsvergelijkingen. 
    
    \begin{question} \( x^2 + 4x + 5    = 0 \) \begin{uitkomst} Er zijn geen oplossingen in de reële getallen.                \end{uitkomst} \end{question}
    \begin{question} \( x^2 + 2x - 8    = 0 \) \begin{uitkomst} Er zijn twee reële oplossingen  \( -4 \) en \( 2 \)           \end{uitkomst} \end{question}
    \begin{question} \( 4x^2 + 12x + 9  = 0 \) \begin{uitkomst} De dubbele wortel is gelijk aan \( -\frac{3}{2} \)            \end{uitkomst} \end{question}
    \begin{question} \( 3x^2 + 6x + 3   = 0 \) \begin{uitkomst} De dubbele wortel is gelijk aan \(  -1 \)                     \end{uitkomst} \end{question}
    \begin{question} \( -x^2 + 4x + 1   = 0 \) \begin{uitkomst} Er zijn twee reële oplossingen  \( -\frac{1}{2} \) en \( 5 \) \end{uitkomst} \end{question}
    
\end{exercise}



\end{document}
